\iffalse
\title{March:2021}
\author{AI24BTECH11008}
\section{mcq-single}
\fi
	\item If P and Q are two statements, then which of the following compound statements is a tautology?\hfill [March 2021]
\begin{enumerate}
    \item [a.] $\brak{\brak{P\rightarrow Q}\land \neg Q}\rightarrow P$
    \item [b.] $\brak{\brak{P\rightarrow Q}\land \neg Q}\rightarrow \neg P$
    \item [c.] $\brak{\brak{P\rightarrow Q}\land \neg Q}$
    \item [d.] $\brak{\brak{P\rightarrow Q}\land \neg Q}\rightarrow \neg Q$
\end{enumerate}
\item  Consider a hyperbola $H : x^2 - 2y^2 = 4$. Let the tangent at a point $P \brak{4, \sqrt{6}}$ meet the x-axis at $Q$ and the latus rectum at $ R \brak{x_1, y_1}$, where  $x_1 > 0$. If $ F $ is a focus of $H$  which is nearer to the point  $P$, then the area of $\Delta QFR$ is equal to:\hfill [March 2021]
\begin{multicols}{2}
\begin{enumerate}
    \item [a.] $\sqrt{6}-1$
    \item [b.] $4\sqrt{6}-1$
	    \columnbreak
    \item [c.] $4\sqrt{6}$
    \item [d.] $\frac{7}{\sqrt{6}}-2$
\end{enumerate}
\end{multicols}
\item Let $f:R\rightarrow R$ be a function defined as 
$$f\brak{x} = \begin{cases}
    \frac{\sin\brak{a+1}x + \sin 2x}{2x}, & x<0\\
    b, & x=0\\
    \frac{\sqrt{x+bx^3-\sqrt{x}}}{bx^{5/2}}, & x>0
\end{cases}$$
If f is continuous at x = 0, then the value of a+b is equal to\hfill [March 2021]
\begin{multicols}{2}
\begin{enumerate}
    \item [a.] -2
    \item [b.] $\frac{-2}{5}$
	    \columnbreak
    \item [c.] $\frac{-3}{2}$
    \item [d.] -3
\end{enumerate}
\end{multicols}
\item Let $y=y\brak{x}$ be the solution of the differential equation $\frac{dy}{dx} =\brak{y+1}\brak{\brak{y+1}e^{x^2/2}-x}$, $0<x<2.1$, with y\brak{2} = 0. Then the value of $\frac{dy}{dx}$ at x=1 is equal to \hfill [March 2021]
\begin{multicols}{2}
\begin{enumerate}
    \item [a.] $\frac{e^{5/2}}{\brak{1+e^2}^2}$
    \item [b.] $\frac{5e^{1/2}}{\brak{e^2}+1}$
	    \columnbreak
    \item [c.] $\frac{-2e^2}{\brak{1+e^2}^2}$
    \item [d.] $\frac{-e^{3/2}}{\brak{e^2+1}^2}$
\end{enumerate}
\end{multicols}
\item Let a tangent be drawn to the ellipse $ \frac{x^2}{27} + y^2 = 1 $ at $ \brak{3\sqrt{3} \cos \theta, \sin \theta}$ where $\theta \in \brak{0, \frac{\pi}{2}}$ . Then the value of $\theta$  such that the sum of intercepts on axes made by a tangent is minimum is equal to: \hfill [March 2021]
\begin{enumerate}
    \item [a.] $\frac{\pi}{8}$
    \item [b.]  $\frac{\pi}{6}$
    \item [c.]  $\frac{\pi}{3}$
    \item [d.]  $\frac{\pi}{4}$
\end{enumerate}

